% Generated draft: requires mathematical and source review.
\documentclass[11pt]{article}
\usepackage[margin=25mm]{geometry}
\usepackage{fontspec}
\setmainfont{Latin Modern Roman}
\usepackage{amsmath,amssymb,mathtools,unicode-math}
\setmathfont{Latin Modern Math}
\usepackage{enumitem,needspace,xurl,hyperref}
\setlength{\parindent}{0pt}
\setlength{\parskip}{5pt}
\setlength{\emergencystretch}{3em}
\newcommand{\N}{\mathbb{N}}
\newcommand{\Z}{\mathbb{Z}}
\newcommand{\Q}{\mathbb{Q}}
\newcommand{\R}{\mathbb{R}}
\newcommand{\C}{\mathbb{C}}
\newcommand{\F}{\mathbb{F}}
\newcommand{\A}{\mathbb{A}}
\newcommand{\PP}{\mathbb{P}}
\newcommand{\E}{\mathbb{E}}
\newcommand{\eps}{\varepsilon}
\newcommand{\dd}{\,\mathrm{d}}
\begin{document}
\section*{UnsolvedMath ID: 20003130}
\textbf{Source problem code:} catalogue entry
\subsection{Definitions and mathematical statement}
A real function on a Euclidean open set is locally DC if locally it is a difference of two convex functions. A map is DC when its coordinate functions are, and a bi-DC homeomorphism has both it and its inverse locally DC. A DC manifold has charts with bi-DC transition maps. Locally finite piecewise-linear (PL) maps are locally DC, and bi-PL homeomorphisms are bi-DC.

The source asks to understand DC structures, whether Perelman's Stability Theorem holds in the DC category, and whether ``PL = DC'' always. The last phrase could concern existence or equivalence of manifold structures rather than literal equality of all maps. A DC stability formulation must also specify whether DC charts are available globally or only on the regular locus of the Alexandrov spaces involved, and the precise noncollapse and curvature hypotheses. These choices are not supplied.

The aid's compactness obstruction is that uniform bilipschitz control does not make the class of bi-DC maps closed under uniform convergence. One concrete accumulating-interval version of the stated obstruction is the following. On $\R$, let $a=1$ outside $(0,1/2)$, and on
\[
I_k=(2^{-k-1},2^{-k}),\qquad k\geq1,
\]
let $a$ alternate between $1$ and $2$. Let $a_N$ agree with $a$ on $I_1,\ldots,I_N$ and equal $1$ on the remaining accumulating intervals; set
\[
f(x)=\int_0^x a(t)\,\dd t,\qquad f_N(x)=\int_0^x a_N(t)\,\dd t.
\]
Then the $f_N$ are increasing bi-PL homeomorphisms $\R\to\R$ and satisfy, uniformly in $N$,
\[
|x-y|\leq|f_N(x)-f_N(y)|\leq2|x-y|,\qquad
\|f_N-f\|_\infty\leq2^{-N-1}.
\]
The limit is bilipschitz, but its weak derivative $a$ has infinite variation on every neighborhood of $0$. By the one-dimensional DC/BV characterization, $f$ is not locally DC at $0$. Taking $(x_1,\ldots,x_d)\mapsto(f_N(x_1),x_2,\ldots,x_d)$ gives the corresponding local obstruction in every dimension; restriction to a line detects the failure of DC in the limit.

The aid's positive one-dimensional repair is: if increasing homeomorphisms $f_j$ of a fixed interval satisfy uniform two-sided bilipschitz bounds, converge uniformly to $f$, and their weak derivatives satisfy a uniform total-variation bound, then $f'$ is of bounded variation and $f$ is bi-DC. Local uniform BV bounds give the local version. The inverse conclusion uses monotonicity and the derivative lower bound; it is not merely a conclusion from uniform convergence alone.

These results distinguish first-order bilipschitz control from the second-order measure control encoded by BV derivatives. They do not refute existence of DC comparison homeomorphisms in a correctly specified Alexandrov stability theorem, nor decide equivalence of PL and DC manifold structures. The maps in the counterexample are a closure test, not purported Perelman stability maps for a prescribed convergent family of spaces.
\subsection{Short English statement}
Understand DC structures and a DC version of stability. The supplied aid shows why uniform bilipschitz limits alone need not remain DC and gives a one-dimensional BV compactness repair, leaving the actual stability and PL/DC category questions unresolved.
\subsection{Sources}
\begin{itemize}
\item Local UnsolvedMath registry, permanent ID $20003130$, source code \texttt{catalogue entry}; canonical \texttt{aim-topology-notes.json}, \texttt{notes[217]}.
\item AIM workshop \emph{Manifolds with nonnegative sectional curvature}, Problem 5: \url{https://aimath.org/WWN/nnsectcurvature/nnsectcurvature.pdf}. The source asks about DC structures, Perelman's Stability Theorem and ``PL = DC''. Original remarks and literature field are empty; ``Sta-bility'' is line-break typography.
\item Research attempt 1: \texttt{partial\_result}, status at run \texttt{open}, current status \texttt{partially\_solved}, classification \texttt{PARTIAL-PROGRESS}, rigor \texttt{proved}, novelty confidence low. No authoritative complete references or full DC-stability formulation are supplied. The machine-generated aid is not peer reviewed and requires independent verification before citation.
\end{itemize}
\end{document}
